\section{问题一的建模与求解}
\label{sec:problem1}
\subsection{优化模型}
\subsubsection{目标函数}
以 $1\,\mathrm{m^3}$ 拌合物为计算基准，建立双目标模型：
\begin{equation}
  \min_{\boldsymbol{x}}\bigl(E(\boldsymbol{x}),C(\boldsymbol{x})\bigr),
  \qquad E(\boldsymbol{x})=\sum_i e_i x_i,\quad
  C(\boldsymbol{x})=\sum_i c_i x_i.
  \label{eq:objectives}
\end{equation}
其中 $E$ 为碳排放，$C$ 为成本。

\subsubsection{约束条件}
可行域需满足体积守恒、强度下限及材料用量边界。
实际计算前，应补充响应函数、约束表达式和参数来源。

\subsection{求解方法}
\subsubsection{算法原理}
采用 NSGA-II 搜索非支配方案\citep{deb2002}，记录种群规模、迭代次数与随机种子。
当两个目标均最小化时，“支配”指各目标均不劣且至少一个严格更优。
\begin{theorem}[支配关系的传递性]
  若方案 $\boldsymbol{x}$ 支配 $\boldsymbol{y}$，且 $\boldsymbol{y}$ 支配
  $\boldsymbol{z}$，则 $\boldsymbol{x}$ 支配 $\boldsymbol{z}$。
\end{theorem}

\subsubsection{求解流程}
基本流程见图~\ref{fig:framework}，三线表伪代码见算法~\ref{alg:nsga2}。
伪代码仅展示求解框架，具体算子与约束处理需按实际问题补充。
\begin{figure}[htbp]
  \centering
  % Vector example; replaced automatically by framework.pdf or framework.png.
\begin{tikzpicture}[
  stage/.style={draw=black!65,rounded corners=2pt,align=center,
    minimum width=38mm,minimum height=13mm,font=\normalsize},
  >={Stealth[length=2mm]}
]
  \node[stage] (data) {数据检查\\范围与异常值};
  \node[stage,right=8mm of data] (model) {模型估计\\最小二乘法};
  \node[stage,right=8mm of model] (check) {结果检验\\误差与残差};
  \draw[->] (data) -- (model);
  \draw[->] (model) -- (check);
\end{tikzpicture}

  \caption{建模与求解流程示例}
  \label{fig:framework}
\end{figure}

% Float the complete narrow algorithm, including its title and three rules.
\begin{figure}[!htb]
  \centering
  \begin{minipage}{0.86\linewidth}
    \setlength{\intextsep}{0pt}
    \begin{algorithm}[H]
      \caption{NSGA-II 求解框架（伪代码示例）}
      \label{alg:nsga2}
      \begin{algorithmic}[1]
        \Require 种群规模 $N$，最大迭代次数 $T$，可行域 $\Omega$
        \Ensure 最终种群的第一非支配前沿
        \State 初始化可行种群 $P_0$，计算目标值 $(E,C)$
        \State 对 $P_0$ 作非支配排序，计算各前沿的拥挤距离
        \For{$t=0,\ldots,T-1$}
          \State 按非支配等级与拥挤距离选择父代，交叉、变异生成 $Q_t$
          \State 修复不可行解，计算 $Q_t$ 的目标值
          \State 对 $P_t\cup Q_t$ 作非支配排序，计算各前沿的拥挤距离
          \State 按等级选入前沿，超额时按拥挤距离降序截取，留 $N$ 个解为 $P_{t+1}$
        \EndFor
        \State \Return $P_T$ 中的第一非支配前沿
      \end{algorithmic}
    \end{algorithm}
  \end{minipage}
\end{figure}
